%==============================================================================
% The Cost of Unlinkability:
% An Empirical Evaluation of A2L Anonymous Atomic Swaps on Bitcoin Taproot
%
% Target venue : ESEM 20XX (short paper, 4-6 pages, English)
% Template     : ACM acmart, sigconf style (ESEM uses ACM proceedings format)
% Compile with : latexmk -pdf paper.tex
%==============================================================================
\documentclass[sigconf,nonacm,natbib=false]{acmart}

% ----- packages --------------------------------------------------------------
\usepackage[T1]{fontenc}
\usepackage[utf8]{inputenc}
\usepackage{booktabs}
\usepackage{siunitx}
\usepackage{microtype}
\usepackage{xspace}
\usepackage{listings}
\usepackage{hyperref}
\usepackage{cleveref}
\usepackage{amsmath}
\usepackage{amssymb}
\usepackage{graphicx}
\usepackage{xcolor}

% ----- macros ----------------------------------------------------------------
\newcommand{\toolname}{\textsc{Tortuga}\xspace}
\newcommand{\atwol}{A\textsuperscript{2}L\xspace}
\newcommand{\rqlabel}[1]{\textbf{RQ#1}}
\newcommand{\hyplabel}[1]{\textbf{H#1}}
\newcommand{\code}[1]{\texttt{\small #1}}

% si-units configuration -- decimal separator: dot (English convention)
\sisetup{detect-all, group-separator={\,}, output-decimal-marker={.}}

% ----- metadata --------------------------------------------------------------
\settopmatter{printacmref=false, printccs=false, printfolios=true}
\setcopyright{none}
\acmConference[ESEM~20XX]{ACM/IEEE Intl.\ Symposium on Empirical Software Engineering and Measurement}{20XX}{TBD}

% ----- title -----------------------------------------------------------------
\title{The Cost of Unlinkability:\\
An Empirical Evaluation of \atwol Anonymous Atomic Swaps on Bitcoin Taproot}

\author{Yan Fernandes}
\affiliation{%
  \institution{UCSal}
  \city{Salvado}
  \country{Brazil}
}
\email{yanferpi@proton.me}

% ----- begin document --------------------------------------------------------
\begin{document}

\begin{abstract}
	Submarine swaps used by production Bitcoin services (Boltz, Loop) rely on
	Hash Time-Locked Contracts (HTLCs) whose preimage hash is identical on both
	legs of a swap, allowing trivial linking of sender and receiver by any
	observer. \emph{Anonymous Atomic Locks} (\atwol, Tairi et al., IEEE S\&P~2021)
	remove this correlation by replacing the shared hash with Schnorr adaptor
	signatures locked to randomizable Castagnos--Laguillaumie (CL) puzzles. While
	the cryptographic security of \atwol is well established, the
	\emph{engineering cost} of adopting it in production-grade Bitcoin software
	remains unquantified. We present a controlled, cross-architecture experiment
	comparing a Rust reference implementation of \atwol on Bitcoin Taproot
	(\toolname) against an HTLC baseline. We measure end-to-end latency,
	on-chain footprint, fees, and peak memory across $N=\num{30}$ runs per arm on
	two heterogeneous machines (Apple~M4~Pro/ARM64 and AMD~EPYC-Genoa/x86\_64).
		[Results placeholder: report effect sizes for H1--H4 with 95\,\% CIs.]
	Our contribution is the first empirical characterisation of the
	unlinkability--cost trade-off for \atwol on Bitcoin and a fully replicable
	benchmarking harness for future privacy-preserving swap research.
\end{abstract}

\keywords{Empirical software engineering, atomic swaps, privacy, Bitcoin,
	adaptor signatures, replication study}

\maketitle

% =============================================================================
\section{Introduction}
\label{sec:intro}

Bitcoin submarine swaps---the operational backbone of off-chain liquidity
services such as Boltz and Lightning Loop---use HTLCs in which the same
SHA-256 preimage hash $H$ appears on both legs of the swap. Although each
leg is a separate transaction, any observer can match the two by inspecting
$H$, breaking unlinkability between payer and payee. The privacy of the user
therefore depends entirely on the trustworthiness of the swap provider.

The Anonymous Atomic Locks construction (\atwol) of \citet{tairi2021a2l}
removes this trust assumption. \atwol replaces the shared hash with a pair
of \emph{randomized adaptor signatures} bound to a CL-encrypted puzzle;
each leg of the swap commits to a different adaptor point on-chain, so the
two transactions are cryptographically unlinkable even by the provider.

While \atwol's cryptographic guarantees have been proven, its
\emph{engineering cost} in a production-grade Bitcoin setting (Taproot,
BIP-340 Schnorr) has not been empirically characterised. Practitioners need
quantitative answers to deploy this privacy improvement.

\paragraph*{Research Questions.}
\begin{itemize}
	\item \rqlabel{1} What is the end-to-end latency overhead of \atwol
	      compared to HTLC submarine swaps?
	\item \rqlabel{2} What is the on-chain footprint difference (vbytes,
	      fees in sat/vB) between \atwol and HTLC?
	\item \rqlabel{3} What is the peak memory and per-phase computational
	      overhead of \atwol relative to HTLC?
\end{itemize}

\paragraph*{Contributions.}
(i)~\toolname, an open-source Rust implementation of \atwol on Bitcoin
Taproot (5~crates, \num{82}~tests).
(ii)~A pre-registered, controlled, cross-architecture experiment
comparing \atwol against an HTLC baseline.
(iii)~A fully replicable benchmarking harness (instrumented CSV emission,
analysis pipeline, Dockerfile, frozen commit) archived on Zenodo.

% =============================================================================
\section{Background}
\label{sec:background}

\subsection{HTLC Submarine Swaps}
A submarine swap exchanges off-chain liquidity (e.g., Lightning) for on-chain
funds. Both legs are funded with the script\linebreak
\code{IF <hash> EQUALVERIFY <pubA> CHECKSIG ELSE <timeout> CLTV DROP <pubB> CHECKSIG ENDIF}.
The recipient learns the preimage $r$ where $\text{SHA256}(r) = H$ to claim
funds; revealing $r$ unlocks the other leg. The shared $H$ makes the two
transactions trivially linkable.

\subsection{Adaptor Signatures and CL Puzzles}
A Schnorr adaptor signature is a pre-signature $\tilde{\sigma}$ over an
adaptor point $T = t \cdot G$; revealing the full signature $\sigma$
discloses $t$. \atwol composes adaptor signatures with a Castagnos--Laguillaumie
public-key encryption scheme whose ciphertexts form a \emph{randomizable
	puzzle}: anyone with the puzzle and a fresh randomizer $\rho$ can produce a
new puzzle $Z'$ for $t + \rho$ without knowledge of $t$. The unlinkability
of \atwol follows from puzzle randomizability and Schnorr adaptor security
\citep{tairi2021a2l, aumayr2021generalized}.

% =============================================================================
\section{Related Work}
\label{sec:related}

\paragraph*{Privacy-preserving payment hubs.} TumbleBit \citep{heilman2017tumblebit}
introduced unlinkable hub payments on Bitcoin using RSA puzzles; \atwol
\citep{tairi2021a2l} reduces the cryptographic and on-chain footprint and is
the first construction compatible with Taproot key-path spends. BlindHub
\citep{qin2023blindhub} extends \atwol with blind signatures.

\paragraph*{Empirical evaluation of cryptographic protocols.}
Performance studies of payment-channel protocols
\citep{harris2020flare,roos2018settling} measure throughput and latency in
simulated networks; we focus instead on \emph{per-swap} cost in a real
Bitcoin regtest network with side-by-side baseline.

\paragraph*{Experimental software engineering on cryptographic systems.}
Replication studies remain rare in the cryptographic-systems space
\citep{wohlin2014guidelines,carver2010replications}; our work follows the
guidelines of \citet{wohlin2012experimentation} for design, randomisation,
and threats-to-validity reporting.

% =============================================================================
\section{Methodology}
\label{sec:method}

We follow a single-factor, two-level, between-subjects controlled experiment.
The protocol is pre-registered: \cref{tab:design} summarises the design and
\cref{tab:vars} the operationalisation.

\subsection{Object of Study}
\toolname is a Rust workspace of five crates (\code{cl-crypto},
\code{adaptor}, \code{protocol}, \code{bitcoin}, \code{cli}) totalling 82
unit and integration tests. It exposes a \code{compare} subcommand that runs
\atwol and HTLC swaps side-by-side against a Nigiri Bitcoin regtest node.

\subsection{Independent Variable}
\textsc{Protocol} $\in \{\atwol, \text{HTLC}\}$.

\subsection{Dependent Variables}
\Cref{tab:vars} lists the \num{12} dependent variables with units and source.

\begin{table}[t]
	\small
	\caption{Operationalisation of dependent variables.}
	\label{tab:vars}
	\begin{tabular}{@{}lll@{}}
		\toprule
		Variable                 & Unit & Source                              \\
		\midrule
		\code{e2e\_latency\_ms}  & ms   & wall-clock init$\to$claim           \\
		\code{pgen\_ms}          & ms   & \code{cl-crypto} (A$^2$L)           \\
		\code{prand\_ms}         & ms   & \code{cl-crypto} (A$^2$L)           \\
		\code{psolve\_ms}        & ms   & \code{cl-crypto} (A$^2$L)           \\
		\code{cldl\_prove\_ms}   & ms   & ZK proof (A$^2$L)                   \\
		\code{cldl\_verify\_ms}  & ms   & ZK proof (A$^2$L)                   \\
		\code{adaptor\_*\_ms}    & ms   & \code{adaptor} sign/verify/complete \\
		\code{htlc\_*\_ms}       & ms   & HTLC build/claim                    \\
		\code{tx\_vbytes}        & vB   & \code{rust-bitcoin} \code{vsize()}  \\
		\code{fee\_sats}         & sat  & vbytes $\times$ \SI{5}{sat/\vbyte}  \\
		\code{peak\_rss\_kb}     & kB   & \code{getrusage} (normalised)       \\
		\code{cpu\_user/sys\_ms} & ms   & \code{getrusage}                    \\
		\bottomrule
	\end{tabular}
\end{table}

\subsection{Hypotheses}
\Cref{tab:hyps} states the four pre-registered hypotheses. We use a
significance level $\alpha = 0.05$ with Bonferroni correction across the
four families ($\alpha' \approx 0.0125$).

\begin{table}[t]
	\small
	\caption{Pre-registered hypotheses. All tests are two-sided unless noted.}
	\label{tab:hyps}
	\begin{tabular}{@{}lp{0.78\columnwidth}@{}}
		\toprule
		ID           & Alternative hypothesis H$_1$                             \\
		\midrule
		\hyplabel{1} & median end-to-end latency of \atwol is greater than HTLC \\
		\hyplabel{2} & median on-chain vbytes of \atwol is smaller than HTLC    \\
		\hyplabel{3} & median fee in satoshis of \atwol is smaller than HTLC    \\
		\hyplabel{4} & median peak RSS of \atwol is greater than HTLC           \\
		\bottomrule
	\end{tabular}
\end{table}

\subsection{Experimental Design}
\begin{table}[t]
	\small
	\caption{Experimental design summary.}
	\label{tab:design}
	\begin{tabular}{@{}ll@{}}
		\toprule
		Element       & Choice                                             \\
		\midrule
		Design        & single-factor, between-subjects                    \\
		Factor        & Protocol \{A$^2$L, HTLC\}                          \\
		Sample size   & $N=30$ per arm (60 total) per machine              \\
		Randomisation & run order randomised; container reset/run          \\
		Warm-up       & 3 discarded runs per arm                           \\
		Hardware (M1) & Apple M4 Pro, 24\,GB, macOS                        \\
		Hardware (M2) & AMD EPYC-Genoa 8\,vCPU, 15\,GB, NixOS 24.11        \\
		Toolchain     & rustc 1.82+ \code{--release}, LTO, \code{mimalloc} \\
		Network       & Nigiri regtest, fresh container/run                \\
		Swap amount   & \num{100000} sat                                   \\
		Fee rate      & \SI{5}{sat/\vbyte}                                 \\
		CL security   & \num{1827}-bit class-group discriminant            \\
		\bottomrule
	\end{tabular}
\end{table}

\subsection{Cross-Architecture Replication}
We replicate the full design on two heterogeneous machines: M1 (Apple~M4~Pro,
ARM64, macOS) and M2 (AMD~EPYC-Genoa, x86\_64, NixOS~24.11, kernel~6.6.94).
This provides an external-validity probe against architecture-specific
SIMD/aes-ni effects and OS-specific scheduler noise. \code{getrusage.ru\_maxrss}
units differ between Linux (kB) and macOS (bytes); the harness normalises
to kB before emission.

\subsection{Statistical Analysis Plan}
For each dependent variable: (i)~Shapiro--Wilk normality test;
(ii)~Levene's test for variance equality; (iii)~if normal and equal variance,
Welch's $t$-test, else Mann--Whitney $U$; (iv)~effect size as Cohen's $d$
(parametric) or Cliff's $\delta$ (non-parametric) with 95\,\% bootstrap CI
($\num{10000}$ resamples); (v)~Bonferroni correction across H1--H4.
All analyses pre-registered in \code{analysis/run.py} (deliverable~\#4).

% =============================================================================
\section{Results}
\label{sec:results}

% --- placeholders auto-generated by analysis/run.py ---
\paragraph*{[Placeholder]} Descriptive statistics (median, IQR, mean$\pm$SD)
per arm per machine. See \code{tables/descriptives.tex}.

\paragraph*{[Placeholder]} Hypothesis-test outcomes for H1--H4 with effect
sizes and 95\,\% CIs. See \code{tables/hypotheses.tex}.

\paragraph*{[Placeholder]} Boxplots of per-phase timings.
See \code{figures/timings.pdf}.

\paragraph*{[Placeholder]} Cross-architecture comparison (M1 vs M2).
See \code{figures/cross\_arch.pdf}.

% =============================================================================
\section{Discussion}
\label{sec:discussion}

% --- to be drafted after results are in ---
\paragraph*{[Placeholder]} Interpretation of effect sizes in operational
terms (e.g., how many additional ms per swap; sat per swap saved by
Taproot keypath spends).

\paragraph*{[Placeholder]} Implications for production swap providers
(Boltz/Loop): is the latency overhead amortisable across a payment hub?
Is the on-chain saving enough to subsidise the CL cost?

% =============================================================================
\section{Threats to Validity}
\label{sec:threats}

We follow the taxonomy of \citet{wohlin2012experimentation}.

\paragraph*{Construct.}
Our metrics (latency, vbytes, fees, RSS) directly operationalise
``engineering cost'', but developer-experience cost (LOC, complexity,
on-call burden) is out of scope for this short paper.

\paragraph*{Internal.}
JIT/warm-cache effects mitigated by 3 warm-up runs per arm; background
load mitigated by a \code{systemd-run} cgroup (M2) and exclusive use (M1);
allocator pinned (\code{mimalloc}); run order randomised within blocks of 10.

\paragraph*{External.}
Generalisation to mainnet block latency is limited by regtest's artificial
block intervals; we report \emph{protocol-internal} latency separately from
on-chain confirmation latency. \toolname is the only \atwol implementation
benchmarked; cross-architecture replication mitigates platform-specific
generalisation concerns.

\paragraph*{Conclusion.}
$N=30$ per arm gives power $\geq 0.80$ for Cliff's $\delta \geq 0.33$
(post-hoc verification). Bonferroni correction is conservative but
appropriate. Replication package frozen at commit \code{<hash>} and
published on Zenodo.

% =============================================================================
\section{Conclusion and Future Work}
\label{sec:conclusion}

We present the first empirical evaluation of \atwol on Bitcoin Taproot and
quantify the unlinkability--cost trade-off relative to HTLC submarine swaps.
	[Results placeholder: 1--2 sentence headline finding once data is in.]
Future work includes (i)~developer-experience metrics on \toolname;
(ii)~end-to-end measurements on mainnet/testnet with realistic block latency;
(iii)~benchmarking alternative CL backends.

% =============================================================================
\paragraph*{Replication Package.}
Source, raw CSVs, analysis scripts, and Dockerfile are available at
\url{https://github.com/yan-pi/tortuga-swap} (tag \code{esem20XX-replication})
and archived on Zenodo (DOI on submission).

\bibliographystyle{ACM-Reference-Format}
\bibliography{references}

\end{document}
